% !PRINTED 348
\weberpage{348}
%<<CONT
entsteht, kann $(1,1)$, $(2,2)$ noch in $(0,2)$, $(0,1)$, aber nicht in
$(2,0)$, $(1,0)$ oder in $(1,2)$, $(2,1)$ übergehen. Also ist auch die Gruppe $S_3 Q$ noch nicht
zweifach transitiv, wohl aber ist sie primitiv. Nehmen wir aber, um $S_2$
zu bilden, noch die Substitution

\[
\begin{pmatrix} 0, & 1 \\ -1, & 0 \end{pmatrix}
\]

hinzu, so geht $(1,1)$, $(2,2)$ in $(0,2)$, $(0,1)$ und in $(1,2)$, $(2,1)$ und
$(2,0)$, $(1,0)$ über. Also ist die Gruppe $S_2 Q$ zweifach transitiv und folglich
auch $S_1 Q$ und $S Q$. Es fangen also erst bei $S_2 Q$ die Tripelgleichungen an.
Die Gleichungen mit engeren Gruppen sind noch keine eigentlichen Tripelgleichungen.

Wir können noch andere in $S Q$ enthaltene lineare Gruppen bilden, z.\,B. wenn wir
mit $S'$ die cyklische Gruppe 3$^{\mathrm{ten}}$ Grades:

\[
S' = \begin{pmatrix} 1, & 0 \\ 0, & 1 \end{pmatrix}, \begin{pmatrix} 1, & 0 \\ 1, & 1 \end{pmatrix}, \begin{pmatrix} 1, & 0 \\ -1, & 1 \end{pmatrix}
\]

bezeichnen, die Gruppe 27$^{\mathrm{sten}}$ Grades $S' Q = Q S'$. Durch die Gruppe $S'$
geht $(1,1)$ $(2,2)$ in $(1,1)$ $(2,2)$, $(1,2)$ $(2,1)$, $(1,0)$ $(2,0)$ über, während
$(0,1)$ $(0,2)$ nicht daraus entsteht. Also ist auch dies nicht die Gruppe einer
Tripelgleichung.

Hier wäre noch ein Feld weitergehender Untersuchungen über diese speciellen
metacyklischen Gleichungen 9$^{\mathrm{ten}}$ Grades und die Darstellung ihrer
Wurzeln nach den Principien von \gesperrt{Kronecker} (Bd. I, achtzehnter Abschnitt).

\begin{center}
\S.\,94.\\[2pt]
\gesperrt{Realitätsverhältnisse der Tripelgleichungen.}
\end{center}

Wenn der Rationalitätsbereich reell ist, so können wir über die Realität der Wurzeln
einer Tripelgleichung 9$^{\mathrm{ten}}$ Grades einige allgemeine Schlüsse machen.
Wenn in diesem Falle nämlich in einem Tripel zwei reelle Wurzeln vorkommen, so muss,
wie aus der Gleichung $x_3 = \Theta (x_1, x_2)$ folgt, auch die dritte Wurzel reell
sein. Wenn also überhaupt eine imaginäre Wurzel vorhanden ist, so muss in jedem der
vier Tripel, dem diese Wurzel angehört, mindestens noch eine zweite imaginäre Wurzel
vorkommen, und folglich ist die Anzahl der imaginären Wurzeln mindestens gleich 5,
oder, da die Zahl der imaginären Wurzeln
%CONT>>
